\setchapterpreamble[u]{\margintoc}
\chapter{向量與線性方程}
\labch{vectors-linear-equations}

{\small
我在準備 KAN 研究時，重新整理了這份線性代數筆記。本章採用
Gilbert Strang 在 \emph{Introduction to Linear Algebra} 和 MIT 18.06
中使用的行與列觀點，並收錄我自己的推導、計算，以及供後續程式實驗
使用的筆記。\par}

\section[\zhHans 线性方程组]{\zhHans 线性方程组与
\texorpdfstring{$A\mathbf{x}=\mathbf{b}$}{Ax=b}}

A system of linear equations can be compressed into the matrix equation

\[
    A\mathbf{x}=\mathbf{b},
\]

where \(A\in\mathbb{R}^{m\times n}\) contains the coefficients,
\(\mathbf{x}\in\mathbb{R}^{n}\) contains the unknowns, and
\(\mathbf{b}\in\mathbb{R}^{m}\) is the target vector. Consider

\[
    A=
    \begin{bmatrix}
        2 & -1 & 0 \\
        -1 & 2 & -1 \\
        0 & -3 & 4
    \end{bmatrix},
    \qquad
    \mathbf{b}=
    \begin{bmatrix}
        0 \\
        -1 \\
        4
    \end{bmatrix}.
\]

Writing the equation row by row gives

\begin{align*}
    2x-y       &= 0, \\
    -x+2y-z    &= -1, \\
    -3y+4z     &= 4.
\end{align*}

The vector

\[
    \mathbf{x}=
    \begin{bmatrix}0&0&1\end{bmatrix}^{\mathsf T}
\]

is a solution because direct substitution returns \(\mathbf{b}\).

\begin{remark}
The notation \(A\mathbf{x}=\mathbf{b}\) is not merely shorter. It exposes
structure that is difficult to see when the equations are treated separately.
Two especially useful interpretations are the \emph{row picture} and the
\emph{column picture}.
\end{remark}

\section{The row picture}

Each row of a three-variable system describes a plane in
\(\mathbb{R}^{3}\). Solving the system means finding a point that lies on all
three planes at once. For the example above, the common point is
\((0,0,1)\).

This geometric interpretation is intuitive in two or three dimensions, but it
does not scale well as a visual method. A system with nine unknowns would
involve hyperplanes in \(\mathbb{R}^{9}\), which cannot be drawn directly.
The algebra, however, remains unchanged.

\marginnote{The row picture emphasises constraints: every row adds one
condition that a candidate solution must satisfy.}

The row picture is therefore most useful as a conceptual bridge. It explains
why a solution must satisfy every equation simultaneously, while elimination
provides a systematic procedure that works in any finite dimension.

\section{The column picture}

Let \(\mathbf{a}_1,\mathbf{a}_2,\mathbf{a}_3\) denote the columns of \(A\).
Then

\[
    A\mathbf{x}
    =x\mathbf{a}_1+y\mathbf{a}_2+z\mathbf{a}_3.
\]

For the running example, the same system becomes

\[
    x\begin{bmatrix}2\\-1\\0\end{bmatrix}
    +y\begin{bmatrix}-1\\2\\-3\end{bmatrix}
    +z\begin{bmatrix}0\\-1\\4\end{bmatrix}
    =\begin{bmatrix}0\\-1\\4\end{bmatrix}.
\]

The solution \((x,y,z)=(0,0,1)\) is immediately visible: the target is the
third column of \(A\). The important idea is broader than this convenient
case.

\begin{definition}
A vector \(\mathbf{b}\) is a \emph{linear combination} of
\(\mathbf{a}_1,\ldots,\mathbf{a}_n\) if there are scalars
\(x_1,\ldots,x_n\) such that
\[
    \mathbf{b}=x_1\mathbf{a}_1+\cdots+x_n\mathbf{a}_n.
\]
The set of every such combination is the \emph{column space} of \(A\),
written \(\operatorname{Col}(A)\).
\end{definition}

Consequently,

\[
    A\mathbf{x}=\mathbf{b}
    \quad\text{has a solution if and only if}\quad
    \mathbf{b}\in\operatorname{Col}(A).
\]

Changing the target vector tests a different combination of the same columns.
For example,

\[
    \begin{bmatrix}1\\1\\-3\end{bmatrix}
    =
    \begin{bmatrix}2\\-1\\0\end{bmatrix}
    +
    \begin{bmatrix}-1\\2\\-3\end{bmatrix},
\]

so the corresponding system has the solution
\(\mathbf{x}=[1,1,0]^{\mathsf T}\).

\begin{remark}
The column picture focuses on construction. Solving the system asks:
\emph{Which combination of the available columns produces the target?}
\end{remark}

\section{Matrix--vector multiplication}

The column viewpoint gives the most useful definition of matrix--vector
multiplication. If

\[
    A=\begin{bmatrix}\vert&&\vert\\
    \mathbf{a}_1&\cdots&\mathbf{a}_n\\
    \vert&&\vert\end{bmatrix}
    \quad\text{and}\quad
    \mathbf{x}=\begin{bmatrix}x_1\\\vdots\\x_n\end{bmatrix},
\]

then

\[
    A\mathbf{x}=x_1\mathbf{a}_1+\cdots+x_n\mathbf{a}_n.
\]

For instance,

\[
    \begin{bmatrix}2&5\\1&3\end{bmatrix}
    \begin{bmatrix}1\\2\end{bmatrix}
    =1\begin{bmatrix}2\\1\end{bmatrix}
    +2\begin{bmatrix}5\\3\end{bmatrix}
    =\begin{bmatrix}12\\7\end{bmatrix}.
\]

The familiar row-by-column rule produces the same result:

\[
    \begin{bmatrix}
        2\cdot1+5\cdot2\\
        1\cdot1+3\cdot2
    \end{bmatrix}
    =\begin{bmatrix}12\\7\end{bmatrix}.
\]

Both calculations are valid, but they answer different questions. The
row-by-column rule is efficient for computing entries; the combination-of-
columns rule explains what the product means geometrically.

The corresponding PyTorch expression is set like the calculation above, with
no separate code panel:

\[
\begin{aligned}
    &\CodeLine{import torch}\\
    &\CodeLine{A = torch.tensor([[2., 5.], [1., 3.]])}\\
    &\CodeLine{x = torch.tensor([1., 2.])}\\
    &\CodeLine{b = A @ x}
\end{aligned}
\]

It returns \(\mathtt{tensor([12.,\ 7.])}\), matching the calculation.

\section{Reachability, independence, and invertibility}

For a square matrix \(A\in\mathbb{R}^{n\times n}\), an important question is
whether \(A\mathbf{x}=\mathbf{b}\) can be solved for \emph{every}
\(\mathbf{b}\in\mathbb{R}^{n}\). In the column picture, this asks whether the
columns of \(A\) fill the entire space:

\[
    \operatorname{Col}(A)=\mathbb{R}^{n}.
\]

If the columns are linearly independent and span \(\mathbb{R}^{n}\), then
every target is reachable by exactly one coefficient vector. The matrix is
invertible, and

\[
    \mathbf{x}=A^{-1}\mathbf{b}.
\]

If the columns are dependent, at least one column contributes no genuinely new
direction. The matrix is singular. Some targets then lie outside the column
space and cannot be reached; other targets may have more than one
representation.

\begin{description}
    \item[Invertible case] Every \(\mathbf{b}\in\mathbb{R}^{n}\) has one and
    only one solution.
    \item[Singular case] Depending on \(\mathbf{b}\), the system has no
    solution or infinitely many solutions.
\end{description}

\section{Connection to KAN research}

Later work with neural networks repeatedly uses vectors for data and features,
matrices for batched computation and linear transformations, and spaces for
describing which outputs a model can represent. KANs change how learnable
transformations are parameterised, but they do not remove the need to reason
about dimensions, transformations, independence, and numerical computation.

The practical lesson from this chapter is therefore simple:

\begin{quote}
Treat \(A\mathbf{x}\) as a combination of the columns of \(A\), while keeping
the row equations available as explicit constraints.
\end{quote}

This dual viewpoint will guide the later notes on elimination, matrix
factorisation, vector spaces, eigenvalues, optimisation, and model
implementation.
